\clearpage
\SourceLeaf{605}
\NoSourceHead
\begin{center}\large\bfseries BEMERKUNGEN ZUR ERSTEN ABTEILUNG\\DES ZEHNTEN BANDES.\end{center}
\begin{center}\rule{.1\linewidth}{.4pt}\end{center}
\medskip
Die vorliegende erste Abteilung des zehnten Bandes von \textsc{Gauss}’ Werken enthält zunächst einige kleinere, bisher in den Werken noch nicht abgedruckte Veröffentlichungen, dann eine Reihe von Nachlaßstücken und Briefstellen aus den verschiedenen Gebieten der Reinen Mathematik, die eine erneute Durchforschung der Handschriften von \textsc{Gauss} noch zu Tage gefördert hat, endlich eine Nachbildung des wissenschaftlichen \emph{Tagebuchs} oder \emph{Notizenjournals} und einen Abdruck des Textes dieser wichtigen Urkunde mit ausführlichen Erläuterungen. Die Nachlaßstücke und Briefstellen sind entsprechend den in den Werken bisher befolgten Grundsätzen nach ihrer Zugehörigkeit zur \emph{Arithmetik}, \emph{Algebra}, \emph{Analysis} und \emph{Geometrie} in Gruppen zusammengefaßt worden; über Einzelheiten unterrichtet am besten das nachstehende \emph{Inhaltsverzeichnis}. Es möge nur noch erwähnt werden, daß, mit Rücksicht auf den Zusammenhang mit den hier zum ersten Male veröffentlichten Stücken, einzelne schon im dritten und achten Bande der Werke ganz oder teilweise abgedruckte Aufzeichnungen noch einmal wiedergegeben worden sind; der Gesamtumfang dieser Wiederholungen beträgt etwa 20 Seiten. Zur Bequemlichkeit für eine spätere Nachprüfung wurden die Stellen, wo sich die Urschriften der hier abgedruckten Nachlaßstücke befinden, durch Angabe der zur Zeit im \textsc{Gauss}archiv angewandten Bezeichnungsweise kenntlich gemacht. Über diese Bezeichnungsweise sei hier — ohne damit der in einem folgenden Bande zu gebenden ausführlichen Beschreibung des Nachlasses vorgreifen zu wollen — das nachstehende bemerkt.

Die kleinen Notizheftchen, die \textsc{Gauss} von 1798 ab benutzt und anfangs als \emph{Schedae} bezeichnet hat, tragen den Leitbuchstaben A und werden durch die kleinen Buchstaben a, b, .\,.\,., n unterschieden; wir haben also die Schedae Aa, Ab, .\,.\,., An. Die steif gebundenen \textls[120]{Handbücher}, deren erstes im Jahre 1800 begonnen wurde, sind mit dem Leitbuchstaben B und weiter mit a, b, .\,.\,., h bezeichnet; die den so gebildeten Zeichen Ba, Bb, .\,.\,. vorgestellten Zahlen 15, 16, .\,.\,. entsprechen der ältern Nummerierung der Nachlaßstücke, wie sie der in dem \emph{Verzeichniss der Handschriften im Preussischen Staate}, I.\ \emph{Hannover}, 3.\ \emph{Göttingen}, 3.\ (Berlin 1894), S.\ 101—113, enthaltene Katalog des \textsc{Gauss}schen Nachlasses von \textsc{Wilhelm Meyer} gibt. Neben den Schedae und Handbüchern kommen für die abgedruckten Nachlaßstücke noch einzelne Druckwerke in Betracht, die \textsc{Gauss} besessen und deren unbedruckte Seiten oder Durchschußblätter er mit wissenschaftlichen Aufzeichnungen beschrieben hat, so namentlich der \textsc{Leiste} (siehe oben, S.\ 78, Fußnote), ferner eine Reihe von losen Zetteln, die mit den Leitbuchstaben E beziehungsweise F versehen sind, jenachdem sie zahlentheoretischen oder analytischen Inhalts sind.

Die wiedergegebenen Briefstellen wurden, so weit uns die Urschriften zugänglich waren, nach diesen abgedruckt oder mit ihnen verglichen. Bei den Nachlaßstücken und Briefen wurde die Schreibung der
\clearpage
\SourceLeaf{606}
\ArithmeticHead{577}{R}{BEMERKUNGEN ZUR ERSTEN ABTEILUNG DES ZEHNTEN BANDES.}
\fancyhead[C]{\scriptsize BEMERKUNGEN ZUR ERSTEN ABTEILUNG DES ZEHNTEN BANDES.}
\noindent Vorlage unverändert beibehalten, dagegen bei allem, was von den Bearbeitern herrührt, durchweg die neuere allgemeine deutsche Rechtschreibung angewandt. Wie in den früheren Bänden sind Einschaltungen der Bearbeiter in den \textsc{Gauss}schen Text in eckige Klammern [\,], Abdrücke von Briefstellen und Mitteilungen, die nicht von \textsc{Gauss} herrühren, in geschweifte Klammern $\bigl\{\bigr\}$ gesetzt.

Die Auswahl der in diesem Bande abgedruckten \textls[120]{arithmetischen} Nachlaßstücke\linebreak erfolgte im Einvernehmen mit P.\ \textsc{Bachmann}, der auch die Erläuterungen zu vielen dieser Stücke verfaßt hat. Die \textls[120]{geometrischen} Stücke hat P.\ \textsc{Stäckel} herausgegeben und erklärt. Die Bearbeitung der übrigen Teile dieses Bandes, wie auch die allgemeine Redaktion lag in den Händen von L.\ \textsc{Schlesinger}, der dabei von P.\ \textsc{Stäckel} und vom Unterzeichneten unterstützt wurde. Bei der Erläuterung einzelner Nachlaßteile und Tagebuchaufzeichnungen haben, wie an den betreffenden Stellen kenntlich gemacht ist, außer den bereits Genannten noch M.\ \textsc{Brendel}, R.\ \textsc{Dedekind} †, L.\ \textsc{Fejér}, R.\ \textsc{Fricke}, A.\ \textsc{Galle}, S.\ \textsc{Gundelfinger} †, J.\ \textsc{Horn}, E.\ \textsc{Landau}, A.\ \textsc{Loewy}, O.\ \textsc{Perron}, K.\ \textsc{Schwering} mitgewirkt.

Die dem Text eingefügten Erläuterungen zu den abgedruckten Stücken werden weiterhin durch die \emph{Aufsätze über Gauss’ wissenschaftliche Tätigkeit auf den einzelnen Gebieten der Mathematik} ergänzt. Diese Aufsätze sollen soweit sie sich auf Reine Mathematik beziehen in der zweiten Abhandlung des vorliegenden zehnten Bandes abgedruckt werden, soweit sie Physik, Astronomie und Geodäsie betreffen, zusammen mit den noch nicht veröffentlichten auf diese Gebiete bezüglichen Stücken des Nachlasses im elften Bande. Vor dem Abdruck in den Werken werden sie in der von \textsc{Brendel}, \textsc{Schlesinger} und mir herausgegebenen Sammlung: \emph{Materialien für eine wissenschaftliche Biographie von Gauss} dem Urteil der Fachgenossen unterbreitet.
\begin{flushright}\textsc{F. Klein.}\end{flushright}
\ClosingRule
